% Authored. Numbers come from generated/macros.tex.

\section{Institutional and Accounting Framework}
\label{sec:framework}

\subsection{The subsectors}

Under ESA 2010 the general government sector S.13 consolidates four subsectors
\citep{eurostat_esa2010}:

\begin{itemize}
\item \textbf{S.1311, central government}, comprising all government units with a
national sphere of competence \emph{with the exception of social security units}
\citep{eurostat_central_glossary};
\item \textbf{S.1312, state government}, which does not apply to Portugal;
\item \textbf{S.1313, local government}, which in the Portuguese case is reported
together with the autonomous regions and appears here as Regional and Local
Government;
\item \textbf{S.1314, social security funds}, comprising all social security units
\emph{regardless of the level of government that operates or manages the scheme}
\citep{eurostat_ssf_glossary}.
\end{itemize}

The two emphasised clauses matter for everything that follows, and they are the
source of a recurring misreading.

\textbf{Central Government in the national accounts is not the State in the
budgetary sense.} The statistical subsector excludes social security units by
definition, and it includes many entities --- public institutes, funds,
reclassified public corporations --- that no reader of a State Budget would think of
as the State. A statement about S.1311 is therefore not a statement about the
\emph{Or\c{c}amento do Estado}, and the two should not be substituted for one
another. Where this paper says Central Government it means S.1311 throughout.

\textbf{The social security subsector is defined by function, not by tier.} S.1314
collects social security units whichever level of government runs them. It is
consequently not a level of administration sitting beneath Central Government; it
is a functionally delimited sector alongside it.

\subsection{The identity}

The balancing item is B.9, net lending (+) or net borrowing (--), and it is the
difference between total revenue and total expenditure for the sector concerned
\citep{eurostat_gov_main}. Consolidation across subsectors means the aggregate is
their sum:
\begin{equation}
\Bgg_t = \Bc_t + \Brl_t + \Bssf_t.
\label{eq:identity}
\end{equation}

Equation~\eqref{eq:identity} is definitional. In published data it closes only to
the rounding of the published series, and we carry that residual as a column rather
than assume it away; over the whole panel its largest absolute value is
\MaxClosureResidual\ million euro, which is the rounding unit of the sources.

Because the subsectors are consolidated, transfers between them net out of the
aggregate but not out of the components. A transfer from Central Government to the
Social Security Funds worsens $\Bc$ and improves $\Bssf$ while leaving $\Bgg$
unchanged. This is central to interpreting the composition: the \emph{location} of
a balance within general government depends on the transfer convention, even though
the aggregate does not. We return to this in Section~\ref{sec:robustness} and
deliberately do not construct a transfer-adjusted balance and call it underlying.

\subsection{Two things the identity cannot do}

It cannot support a counterfactual. Writing
$\Bnonssf_t = \Bc_t + \Brl_t$ and observing that $\Bssf_t > |\Bnonssf_t|$ in a
given year says that two recorded numbers stand in a particular ratio. It does not
say what the aggregate balance would have been had the Social Security Funds been
organised differently, because the counterfactual would change contribution rates,
entitlements and transfers simultaneously, and the identity contains no model of
any of them.

It cannot support an attribution of responsibility. A subsector that accounts
arithmetically for most of an annual improvement has not been shown to have
produced it. The decomposition is silent on intent, on policy and on causation, and
we use the word contribution in its arithmetic sense only.
